\section{Conclusion}
\label{sec:conclusion}

This survey studied Physical AI through the lens of LLM-based world knowledge. 
Rather than treating Physical AI as a purely robotics-centric, vision-centric, or cyber-physical topic, we organized recent progress as a roadmap from language-derived world knowledge to multimodal grounding, action grounding, world modeling, policy learning, and embodied agents. 
This perspective highlights the complementary roles of LLMs and world models: LLMs provide semantic, commonsense, procedural, and causal priors, while world models provide predictive and simulative mechanisms for physical dynamics. 
We hope this roadmap clarifies how progress in language modeling, multimodal reasoning, VLA models, world models, and embodied agents can be connected toward more general Physical AI systems.



\section{Limitations}
Limitations
Limitations
Authors are required to discuss the limitations of their work in a dedicated section titled "Limitations" (not counting towards page limit). Papers without this section will be desk rejected. Please confirm that your paper has a limitations section by checking this box.

This question and those that follow are from the Responsible Research Checklist, please see this page for advice on filling it in: https://aclrollingreview.org/responsibleNLPresearch/. Please note that inappropriate or missing answers to checklist questions can be grounds for DESK REJECTION. If your answer to a given question is 'yes' or 'no', rather than 'n/a', the 'elaboration' fields MUST be filled in.
Limitations
We discuss potential risks in the Limitations section, including the risks of over-relying on LLM-derived world knowledge, grounding failures in physical or embodied settings, unsafe deployment of Physical AI systems, and the limitations of current closed-loop evaluation.

\section{LLM Usage}
Assistants